\section{The fiber Legendre transform and its curvature}\label{sec:fiber}

The analytic criterion leaves us a geometric task: construct the base
potential while retaining explicit curvature margins for later changes.
We first derive the metric and its six curvature blocks for a fixed fiber
center. The partial Legendre transform uses, up to our fixed metric
normalization, the fiber momentum of the circle-invariant constructions
of Calabi and Hwang--Singer \citep{Calabi1979,HwangSinger2002}.
We allow the radial profile to vary over the spatial base, which adds
spatial-derivative terms to the curvature. We compute the full tensor
before estimating these additional terms.

All spatial derivatives in this section are taken with the fiber momentum
\(\tau\) fixed, unless a different convention is stated explicitly.
The spatial variables are \(s=(s_1,\ldots,s_n)\), and the fiber variable is
\(y\in\mathbb C\).  Its center \(\zeta\) is constant in this section.
Let \(q(x,s)>0\), where \(x=\log\tau\), and choose \(F\) by
\begin{equation}\label{fiber:F}
 F_x=q^{-1},\qquad \lim_{x\to-\infty}(F(x,s)-x)=0.
\end{equation}
We assume that this normalization is possible, that \(F\to+\infty\)
as \(x\to+\infty\), and that \(q\), regarded as a function of \(\tau\),
is smooth at zero with \(q(0,s)=1\).  All these statements are local
uniformly in \(s\).  For smooth real functions \(B,\psi\), set
\begin{align}
 U(s,\tau)&=B(s)+\tau\psi(s)-\int_0^\tau F(\log v,s)\,dv,
 \label{fiber:U}\\
 \rho&=\log|y-\zeta|^2=F(\log\tau,s)-\psi(s),
 \label{fiber:inverse}\\
 \ell_{\mathrm{rad}}&=U+\tau\rho
 =B+\int_0^\tau q(\log v,s)^{-1}\,dv.
 \label{fiber:potential}
\end{align}
The last identity follows by integration by parts; the boundary term
\(vF(\log v,s)\) tends to zero as \(v\downarrow0\).
Since \(F_x>0\) and its endpoint limits are infinite with the indicated
signs, \eqref{fiber:inverse} determines a unique \(\tau>0\) for every
\(y\ne\zeta\).

Hereafter \(C=-R\) denotes the negative Hermitian curvature tensor:
\begin{equation}\label{fiber:curvature-convention}
 C_{\alpha\bar\beta\gamma\bar\delta}
 =\partial_\alpha\partial_{\bar\beta}g_{\gamma\bar\delta}
 -g^{\mu\bar\nu}
   (\partial_\alpha g_{\gamma\bar\nu})
   (\partial_{\bar\beta}g_{\mu\bar\delta}).
\end{equation}
The sectional-curvature sign associated with this convention is verified
below; in particular, positivity of the relevant contractions of \(C\)
means nonpositive Riemannian sectional curvature.

\begin{lemma}[Metric and regularity of the Legendre transform]
\label{fiber:metric}
Write
\[
 P=\tau q,\qquad V=(B_{i\bar j}),\qquad
 G=(\psi_{i\bar j}),\qquad h=(U_{i\bar j}),\qquad
 \mathcal N=h_\tau.
\]
Where \(h>0\), the complex Hessian of \(\ell_{\mathrm{rad}}\) is
\begin{equation}\label{fiber:metric-form}
 h_{i\bar j}\,ds_i\,d\bar s_j+
 P\left|d\log(y-\zeta)+U_{\tau i}\,ds_i\right|^2.
\end{equation}
It is positive off the axis.  Its spatial Schur complement is exactly
\(h\).  The potential and its complex Hessian extend smoothly across
\(y=\zeta\); on that axis the Hessian in the original splitting is
\(V\oplus e^\psi\), and is positive whenever \(V>0\).
\end{lemma}

\begin{proof}
Use the holomorphic coordinate \(z_0=\log(y-\zeta)\), so that
\(\rho=z_0+\bar z_0\).  Differentiating \(U_\tau+\rho=0\) gives
\[
 U_{\tau\tau}=-P^{-1},\qquad
 \partial_{z_0}\tau=P,\qquad
 \partial_{s_i}\tau=P U_{\tau i}.
\]
Differentiation of \(\ell=U+\tau\rho\) now gives
\[
 g_{0\bar0}=P,\qquad g_{i\bar0}=P U_{\tau i},\qquad
 g_{i\bar j}=h_{i\bar j}+P U_{\tau i}U_{\tau\bar j},
\]
which is \eqref{fiber:metric-form}.  For regularity, write
\(F=\log\tau+f(\tau,s)\).  The equation
\(f_\tau=(q^{-1}-1)/\tau\), with \(f(0,s)=0\), shows that \(f\)
is smooth.  Equation \eqref{fiber:inverse} becomes
\(\tau e^{f(\tau,s)}=e^\psi|y-\zeta|^2\).  Its derivative with respect
to \(\tau\) at zero is one, so the smooth implicit-function theorem
applies.  In particular,
\(\tau=e^\psi|y-\zeta|^2+O(|y-\zeta|^4)\).
Substitution in \eqref{fiber:potential} proves the stated extension and
axis metric.
\end{proof}

We use a local holomorphic change \(z_0\mapsto z_0+f(s)\) to make
\(U_{\tau i}=0\) at a specified point.  The corresponding change of
\(U\) is linear in \(\tau\) with a pluriharmonic spatial coefficient;
it leaves \(h\) and \(\mathcal N\) unchanged.  The quadratic jet of
\(f\) can also make \(U_{\tau ik}=0\) there.  We call these the tilted
logarithmic coordinates.  The coordinate changes used to evaluate a jet
are held fixed while that jet is differentiated.

\begin{lemma}[The six curvature blocks]\label{fiber:blocks}
In tilted logarithmic coordinates, the tensor
\eqref{fiber:curvature-convention} has the following blocks:
\begin{align}
 A:=C_{0\bar0 0\bar0}
   &=P^2P_{\tau\tau}-|\partial P|_h^2,
 \label{fiber:blockA}\\
 A_i:=C_{i\bar0 0\bar0}
   &=P\bigl(\partial_iP_\tau-
       (\mathcal N h^{-1}\partial P)_i\bigr),
 \label{fiber:blockAi}\\
 W_{i\bar j}:=C_{i\bar j0\bar0}
   &=P\bigl(\partial_i\partial_{\bar j}\log P
       +P_\tau\mathcal N_{i\bar j}
       -P(\mathcal N h^{-1}\mathcal N)_{i\bar j}\bigr),
 \label{fiber:blockW}\\
 C_{i\bar0 k\bar0}&=(\nabla_h^{2,0}P)_{ik},
 \label{fiber:blockspin}\\
 K_{ik\bar j}:=C_{i\bar j k\bar0}
   &=P\nabla_i^h\mathcal N_{k\bar j}
      +P_i\mathcal N_{k\bar j}+P_k\mathcal N_{i\bar j},
 \label{fiber:blockK}\\
 L_{i\bar j k\bar l}:=C_{i\bar j k\bar l}
   &=C^h_{i\bar j k\bar l}
      +P(\mathcal N_{i\bar j}\mathcal N_{k\bar l}
          +\mathcal N_{i\bar l}\mathcal N_{k\bar j}).
 \label{fiber:blockL}
\end{align}
Here \(C^h\) and \(\nabla^h\) are the curvature with sign
\eqref{fiber:curvature-convention} and the K\"ahler connection of the
fixed-\(\tau\) spatial metric.  Index positions in the products use the
indicated Hermitian metric, with the usual conjugations.
\end{lemma}

\begin{proof}
In addition to the tilt, first choose spatial coordinates normal for
\(h\) at the point.  On functions of \((s,\tau)\), physical coordinate
derivatives act as
\[
 D_i=\partial_i+P U_{\tau i}\partial_\tau,
 \qquad D_0=P\partial_\tau,
\]
and their conjugates.  At the point the complete list of cubic entries
needed below is
\begin{equation}\label{fiber:cubics}
 \ell_{00\bar0}=PP_\tau,\quad
 \ell_{i0\bar0}=P_i,\quad
 \ell_{00\bar j}=P_{\bar j},\quad
 \ell_{i0\bar j}=P\mathcal N_{i\bar j},\quad
 \ell_{ik\bar0}=0,\quad \ell_{ik\bar j}=0.
\end{equation}
For example, the identity
\(P^2\partial_\tau U_{\tau i}=P_i\), obtained by spatial
differentiation of \(U_{\tau\tau}=-1/P\), and its conjugate give the
second and third entries.
Directly differentiating the metric entries before evaluating gives
\begin{align*}
 D_0D_{\bar0}P&=P(P_\tau^2+PP_{\tau\tau}),\\
 D_iD_{\bar0}P&=P_iP_\tau+P\partial_iP_\tau,\\
 D_iD_{\bar j}P&=P_{i\bar j}+P\mathcal N_{i\bar j}P_\tau,\\
 D_iD_{\bar j}(PU_{\tau k})
   &=P_i\mathcal N_{k\bar j}+P_k\mathcal N_{i\bar j}
      +P\partial_i\mathcal N_{k\bar j},\\
 D_iD_{\bar j}(h_{k\bar l}+PU_{\tau k}U_{\tau\bar l})
   &=\partial_i\partial_{\bar j}h_{k\bar l}
      +P\mathcal N_{i\bar j}\mathcal N_{k\bar l}
      +P\mathcal N_{i\bar l}\mathcal N_{k\bar j}.
\end{align*}
For completeness, the fourth derivative giving the spin block has a
cancellation which is useful later.  Put \(a_i=U_{\tau i}\).  Before
evaluation,
\[
 \ell_{ik\bar0}=Pa_{ik}+P_i a_k+P_k a_i+PP_\tau a_i a_k.
\]
At the point \(a_i=a_{ik}=0\), while
\[
 (a_i)_\tau=P_i/P^2,\qquad
 (a_{ik})_\tau=P_{ik}/P^2-2P_iP_k/P^3.
\]
Applying \(D_{\bar0}=P\partial_\tau\) therefore gives \(P_{ik}\):
the two positive terms \(P_iP_k/P\) cancel the negative doubled term.

Subtract the cubic contractions in
\eqref{fiber:curvature-convention}.  For A the vertical contraction is
\(PP_\tau^2\) and the spatial one is \(|\partial P|_h^2\).
For \(A_i\) they are \(P_iP_\tau\) and
\(P(\mathcal Nh^{-1}\partial P)_i\).  For \(W\) they are
\(P_iP_{\bar j}/P\) and \(P^2(\mathcal Nh^{-1}\mathcal N)_{i\bar j}\).
The first combines with \(P_{i\bar j}\) to give
\(P\partial_i\partial_{\bar j}\log P\).
For the purely spatial block the contraction is the one defining
\(C^h\).  These calculations prove all formulas in normal spatial
coordinates.  Without spatial normalization, the last entry of
\eqref{fiber:cubics} is \(\partial_i h_{k\bar j}\); its contractions
replace \(P_{ik}\) by \(\nabla^h_i\partial_kP\) and
\(\partial_i\mathcal N_{k\bar j}\) by
\(\nabla_i^h\mathcal N_{k\bar j}\).  This also verifies the displayed
tensorial formulas in general spatial coordinates.
\end{proof}

The six blocks give the complete tensor, but the desired sign concerns
real sectional curvature. The next lemma identifies the matrix tests
corresponding to real two-planes and records exactly how much of the
diagonal curvature can absorb each mixed block.

We specify the norm convention needed to keep all subsequent constants
independent of \(n\).  The norm of a covariant tensor is its
multilinear operator norm on unit vectors in the stated fixed frame.
The norm of a Hermitian test matrix is its Frobenius norm.  Thus, for a
Hermitian matrix \(T\) of rank at most two and a covariant Hermitian form \(Q\),
\begin{equation}\label{fiber:ranknorm}
 |Q:T|\le\sqrt2\|Q\|\,|T|.
\end{equation}
Indeed diagonalize \(T\); the sum of the absolute values of its at most two
nonzero eigenvalues is at most \(\sqrt2|T|\).
The same decomposition bounds a curvature pairing \(C(T,T)\) by
\(2\|C\||T|^2\).  Contractions of cubic tensors through an inverse
metric can instead be regarded as inner products of two covectors.
Their operator bounds therefore involve no summation factor depending
on the dimension.

\begin{lemma}[Real planes and rank-two absorption]\label{fiber:ranktwo}
If \(C(H,H)\ge0\) for every Hermitian contravariant matrix \(H\) of rank
at most two, the underlying real sectional curvatures are nonpositive.
Write the spatial, cross, and fiber diagonal blocks of \(H\) as
\(T,b_\times,t\).  The contraction has the terms
\begin{equation}\label{fiber:contraction}
 At^2+2t(W:T)+2W(b_\times,\bar b_\times)+L(T,T),
\end{equation}
and its other terms have absolute value at most
\begin{equation}\label{fiber:otherterms}
 2|\nabla_h^{2,0}P(b_\times,b_\times)|
 +4|tA_i b_\times^i|+4|(K:T)b_\times|.
\end{equation}
In particular the following quantitative criterion suffices.  Suppose
\(A>0\), \(W>0\), and a nonnegative quadratic form \(\mathcal R(T)\)
satisfies
\begin{align}
 L(T,T)-\frac{|W:T|^2}{.98A}&\ge\mathcal R(T),
 \label{fiber:reserve}\\
 |\nabla_h^{2,0}P(b,b)|&\le\tfrac14 W(b,\bar b),
 \label{fiber:spinabs}\\
 |A_i b^i|&\le\eta\sqrt{A W(b,\bar b)},
 \label{fiber:Aiabs}\\
 |(K:T)b|&\le\eta\sqrt{\mathcal R(T)W(b,\bar b)},
 \label{fiber:Kabs}
\end{align}
for all such tests, where \(\eta\le1/100\).  Then
\begin{equation}\label{fiber:retained}
 C(H,H)\ge .01At^2+W(b_\times,\bar b_\times)
                         +\tfrac12\mathcal R(T).
\end{equation}
All constants are independent of dimension.
\end{lemma}

\begin{proof}
For real vectors with \((1,0)\) parts \(v,w\), set
\(H=i(v\otimes\bar w-w\otimes\bar v)\).  This is Hermitian of
rank at most two.  K\"ahler symmetry gives
\[
 C(H,H)=2\{C(v,\bar v,w,\bar w)
                   -\operatorname{Re}C(v,\bar w,v,\bar w)\}.
\]
Expand \(R(v+\bar v,w+\bar w,w+\bar w,v+\bar v)\).  The only four
nonzero terms are the two alternating type patterns and their
conjugates.  In the convention in which a real vector with complex
components \(v\) has squared length \(g_{i\bar j}v^i\bar v^j\), the complex
bilinear extension of the real metric on a holomorphic and an
antiholomorphic vector is one half the Hermitian pairing.  Accordingly
the real curvature numerator is \(-\tfrac12 C(H,H)\).  This proves the
sign assertion.
Expanding a general block matrix \(H\) gives
\eqref{fiber:contraction}--\eqref{fiber:otherterms}; each mixed
conjugate pair accounts for its factor two or four.  A principal block
\(T\) has rank at most two.

For the quantitative assertion use
\(2t(W:T)\ge-.98At^2-|W:T|^2/(.98A)\).
The remaining diagonal coefficient is .02A.  The spin term costs at
most \(W/2\).  The other two terms obey
\[
 4\eta |t|\sqrt{AW}\le .01At^2+400\eta^2W,
 \qquad
 4\eta\sqrt{\mathcal R W}\le\tfrac12\mathcal R+8\eta^2W.
\]
Since \(1.5-408\eta^2\ge1\) when \(\eta\le1/100\),
\eqref{fiber:retained} follows.
\end{proof}

\begin{lemma}[Exact perturbation identity]\label{fiber:perturbation}
In any fixed holomorphic coordinates let \(H\) denote the cubic jet viewed
as a map from two tangent inputs to a covector, and write \(H=\Gamma g\).
For a perturbation with \(g+\Delta g>0\), the change of the cubic
contraction \(Hg^{-1}H^*\) is exactly
\begin{equation}\label{fiber:perturbformula}
 \Gamma\Delta H^*+\Delta H\Gamma^*-\Gamma\Delta g\Gamma^*
 +(\Delta H-\Gamma\Delta g)(g+\Delta g)^{-1}
                 (\Delta H-\Gamma\Delta g)^*.
\end{equation}
Consequently, if \(\|g^{-1}\Delta g\|\) in metric-unit axes is at most
one half, the inverse of the perturbed metric is at most twice that
of \(g\), and
\begin{align}\label{fiber:perturbbound}
 \|\Delta C\|\le {}&\|\Delta j^4\ell\|
       +2\|\Gamma\|\|\Delta H\|
       +\|\Gamma\|^2\|\Delta g\|\\
 &+2\|g^{-1}\|
       (\|\Delta H\|+\|\Gamma\|\|\Delta g\|)^2.
 \notag
\end{align}
The same bound, with an absolute factor, controls rank-two contractions.
\end{lemma}

\begin{proof}
Set \(Q=\Delta H-\Gamma\Delta g\).  Then
\(H+\Delta H=\Gamma(g+\Delta g)+Q\).  Multiply out the corresponding
cubic contraction and subtract \(\Gamma g\Gamma^*\).  The two cross
terms and \(\Gamma\Delta g\Gamma^*\) combine to give the first three
terms in \eqref{fiber:perturbformula}.  Taking operator norms gives
\eqref{fiber:perturbbound}.  The inverse comparison is the elementary
spectral estimate for \(I+g^{-1/2}\Delta g g^{-1/2}\); the final statement
uses \eqref{fiber:ranknorm} twice.
\end{proof}
